\documentclass[10pt,twocolumn]{article}

\usepackage[margin=0.75in]{geometry}
\usepackage{amsmath}
\usepackage{amssymb}
\usepackage{booktabs}
\usepackage{array}
\usepackage{tabularx}
\usepackage{times}
\usepackage{url}

\pagestyle{plain}
\providecommand{\keywords}[1]{\par\smallskip\noindent\textbf{Keywords: }#1\par}
\newcolumntype{Y}{>{\raggedright\arraybackslash}X}

\title{Adversarially Robust and Controllable Execution of Ontology Workflows}

\author{Anonymous Authors\\Anonymous Institution}

\begin{document}
\maketitle

\begin{abstract}
Ontology Workflow Skills provide typed operations, object bindings, state predicates, and control-flow constraints, but their execution by LLM agents remains vulnerable to benign faults and adversarial manipulation. Tool outputs may be incomplete or malicious, tool descriptions may hijack selection, user instructions may induce policy violations, and runtime state may be stale or poisoned. We formulate robust controlled workflow execution: preserve task utility and workflow constraints under bounded noise and attack while preventing irreversible unsafe effects. We propose an Ontology-Compiled Execution Envelope that separates trusted workflow semantics from untrusted natural-language content, maintains uncertainty over runtime ontology state, restricts candidate actions to workflow-consistent object and state transitions, applies worst-case admissibility to irreversible operations, and permits bounded information acquisition and recovery for reversible operations. Evaluation reports clean utility, utility under attack, attack success, constraint violations, unsafe effects, over-blocking, degradation slope, robust radius, and empirical worst-case lower bounds.
\end{abstract}

\keywords{LLM agents, robust execution, ontology workflows, prompt injection, controllability}

\section{Introduction}
LLM agents increasingly execute long-horizon tasks through external tools. Their operational environment is neither clean nor trustworthy. Tool calls can time out, responses can omit fields, schemas can drift, user goals can change, and external content can contain malicious instructions. These failures are especially consequential when agents execute workflows with irreversible side effects.

Prompt-injection research shows that tool-integrated agents can be hijacked by instructions embedded in external content. InjecAgent reports substantial indirect-prompt-injection vulnerability across tool agents~\cite{zhan2024injecagent}. AgentDojo evaluates utility and security in dynamic stateful environments~\cite{debenedetti2024agentdojo}. Agent Security Bench covers prompt injection, memory poisoning, backdoors, and mixed attacks while introducing utility-security tradeoff metrics~\cite{zhang2024asb}. Adaptive attacks further show that defenses evaluated only against fixed attacks are unreliable~\cite{zhan2025adaptive}.

Runtime constraint systems provide another line of defense. AgentSpec enforces custom rules, Progent controls privileges, and Agent-C enforces temporal properties~\cite{wang2025agentspec,shi2025progent,kamath2025agentc}. These methods establish that agent behavior can be restricted, but strict restriction can reduce task success. A robust method must therefore improve the utility--control boundary, not merely block more actions.

We study this problem for agents executing Ontology Workflow Skills. Such workflows define allowed function nodes, ontology object variables, typed parameter bindings, state predicates, effects, invariants, and exceptional branches. We compile this structure into an Execution Envelope that bounds valid runtime transitions. The paper does not claim universal formal safety. Natural-language attack spaces are effectively unbounded. Instead, we define an explicit threat model, attack budget, and adaptive adversary, then report empirical worst-case lower bounds under that model.

The contributions are:
\begin{itemize}
    \item We define a threat model for ontology workflow execution covering benign faults, indirect injection, tool hijacking, state poisoning, and objective drift.
    \item We propose an Ontology-Compiled Execution Envelope with typed trust boundaries, uncertainty-aware state, risk-sensitive admissibility, and bounded recovery.
    \item We design an adaptive evaluation protocol that emphasizes utility--control frontiers, worst-group performance, and empirical worst-case behavior.
\end{itemize}

\section{Related Work}
\subsection{Tool-Agent Robustness and Security Benchmarks}
InjecAgent targets indirect prompt injection in tool-integrated agents~\cite{zhan2024injecagent}. AgentDojo provides stateful tasks, attacks, and defenses over untrusted tool data~\cite{debenedetti2024agentdojo}. Agent Security Bench formalizes attacks at multiple agent stages~\cite{zhang2024asb}. MCP Security Bench and AgentLAB extend attacks to tool discovery, invocation, response handling, tool chaining, objective drift, and memory poisoning~\cite{zhang2025mcpsecurity,jiang2026agentlab}. These benchmarks motivate our threat model. Our method differs by using ontology workflow semantics to define allowed transitions rather than relying only on text filtering or attack detection.

\subsection{Prompt Injection and Tool Hijacking}
Adaptive attacks bypass multiple prompt-injection defenses~\cite{zhan2025adaptive}. ToolHijacker manipulates tool retrieval and selection through malicious tool descriptions~\cite{shi2025toolhijacker}. Tool-result parsing and trusted/untrusted content separation reduce instruction leakage~\cite{beurerkellner2025patterns,yu2026toolparsing}. We adopt this separation but add ontology object, state, effect, and workflow constraints. A syntactically clean response can still induce an invalid state transition.

\subsection{Runtime Enforcement}
AgentSpec, Progent, solver-aided policy verification, and Agent-C enforce runtime rules, privileges, or temporal constraints~\cite{wang2025agentspec,shi2025progent,winston2026solver,kamath2025agentc}. Verifier Tax shows that long-horizon agents may fail after verifier intervention~\cite{sah2026verifiertax}. This motivates measuring over-blocking and repair outcomes rather than treating every rejection as success.

\section{Method}
\subsection{Threat Model}
The agent executes a gold or extracted Ontology Workflow $W$ through a registered Function Layer $F$. The adversary can perturb untrusted channels within budget $B$: user messages, retrieved text, free-text tool results, tool names or descriptions, non-authoritative memory, and selected state observations. The adversary cannot modify the signed workflow, Function contracts, trusted validators, or benchmark environment state directly. Benign faults share some channels but are sampled independently rather than optimized adversarially.

\subsection{Execution Envelope}
Workflow $W$ is compiled into:
\[
\mathcal{E}=\langle C,A,V,R,P,T,I,\beta\rangle,
\]
where $C$ are control states, $A$ are enabled functions, $V$ are typed object variables, $R$ are admissible bindings, $P$ are state predicates, $T$ is the transition relation, $I$ are invariants, and $\beta$ is the side-effect and recovery budget. A proposed function is accepted only if it lies inside this envelope for the current control and ontology state.

\subsection{Trusted Control Plane and Tainted Data Plane}
Signed workflows and Function contracts are trusted control data. User text, external documents, tool descriptions, and free-text tool outputs are untrusted data unless converted into typed observations through validated parsers. Untrusted text can supply values, but it cannot authorize a new workflow transition, add a new function, or override a policy invariant.

\subsection{Uncertainty-Aware Ontology State}
The runtime maintains an uncertainty set:
\[
\mathcal{U}_t=\{s: s \text{ is consistent with trusted observations up to } t\}.
\]
Missing, stale, and conflicting values widen $\mathcal{U}_t$. Untrusted text does not reduce uncertainty. Structured observations narrow the set only after datatype, provenance, object identity, and temporal consistency checks.

\subsection{Risk-Sensitive Admissibility}
Functions are divided into information or reversible operations and irreversible operations such as mutation, transfer, deletion, payment, disclosure, or committed external effects. For irreversible action $a$, the system requires:
\[
\forall s\in\mathcal{U}_t,\quad \mathrm{Precondition}(a,s)\wedge \mathrm{Invariant}(a,s).
\]
Information and reversible actions may be used for bounded exploration when they reduce uncertainty and remain inside the workflow. This avoids the failure mode in which conservative blocking prevents the agent from acquiring the state needed for a safe decision.

\subsection{Bounded Recovery}
When a proposed action is rejected or a benign fault occurs, recovery is limited to workflow-authorized actions: acquire a declared observation, retry an idempotent function, rebind to a validated object, select an allowed branch, execute registered compensation, request missing user input, or escalate. The LLM can rank recovery candidates but cannot introduce an operation outside the Function Layer or workflow envelope.

\subsection{Adaptive Boundary Search}
An attack generator searches perturbations that maximize unsafe effects, workflow violations, or task failure under budget $B$. It mutates tool text, user instructions, registry metadata, observation values, and multi-turn timing. Population-based search retains successful attacks across rounds. This evaluates the boundary of the method rather than only fixed attack templates.

\section{Experiments}
\subsection{Environments}
Experiments use AgentDojo for indirect prompt injection and utility under attack~\cite{debenedetti2024agentdojo}, InjecAgent for tool-integrated attack cases~\cite{zhan2024injecagent}, Agent Security Bench for mixed attack types~\cite{zhang2024asb}, tau-bench Retail and Airline for stateful policy workflows~\cite{yao2024taubench}, SOP-Bench for long-horizon SOP execution~\cite{nandi2025sopbench}, and controlled workflow microbenchmarks with exact attack targets and state transitions. Primary experiments use gold workflows and Function contracts; extracted artifacts are evaluated separately as a cascade setting.

\subsection{Perturbations}
Benign noise includes timeout, rate limit, partial responses, duplicated responses, missing or renamed fields, stale and conflicting state, similar tools, benign paraphrase, and incomplete user input. Adversarial attacks include indirect injection, malicious tool descriptions, tool-selection hijacking, out-of-scope parameter requests, memory or state poisoning, objective drift, false-error escalation, and mixed multi-turn attacks.

\subsection{Baselines and Ablations}
Baselines include ReAct, workflow text in prompt, instruction hierarchy or text-only isolation, tool-result parsing, prompt-injection filters, action-local schema verification, AgentSpec, Agent-C where reproducible, Execution Envelope without uncertainty sets, Execution Envelope without workflow control, Execution Envelope without adaptive hardening, and the full method.

\subsection{Metrics}
We report clean task success, utility under attack, benign-noise success, attack success rate, workflow conformance rate, constraint violation rate, unsafe side-effect rate, object-binding violation rate, over-blocking rate, adversarial regret, degradation slope over attack budget, worst-attack success, worst-group success across attack families and domains, 95\% lower confidence bound of robust success, robust radius, Net Resilient Performance, recovery success, unnecessary recovery, repair steps, escalation rate, latency, token cost, and tool-call overhead.

\subsection{Empirical Worst-Case Protocol}
For task $i$ and allowed attack set $\mathcal{A}(B)$:
\[
\widehat{WCS}(B)=\frac{1}{N}\sum_i \min_{a\in\mathcal{A}(B)}\mathrm{Success}(i,a).
\]
Because $\mathcal{A}(B)$ is finite and attack search is incomplete, this is an empirical lower bound under the evaluated threat model, not a formal global guarantee. We report bootstrap and binomial lower confidence bounds.

\section{Analysis}
The full method must improve the utility--control frontier. Evidence is sufficient only if it lowers attack success at matched clean utility, or raises utility under attack at matched violation rate, and improves worst-group or empirical worst-case success. More rejected actions alone are not evidence of robustness. We separately analyze clean-utility tax, over-blocking, attack-family attribution, worst-group failures, and whether a new domain can be handled by replacing ontology and workflow definitions without domain-specific control code.

\section{Conclusion}
This paper reframes ontology-grounded workflow execution as a robustness and controllability problem. The Execution Envelope uses workflow, object, state, and effect semantics to constrain how untrusted information can influence actions. Its value must be demonstrated through utility under attack, constraint violations, unsafe effects, degradation curves, robust radius, and empirical worst-case lower bounds.

\begin{thebibliography}{99}

\bibitem{zhan2024injecagent}
Q. Zhan et al.
\newblock InjecAgent: Benchmarking Indirect Prompt Injection in Tool-Integrated Agents.
\newblock arXiv:2403.02691, 2024.

\bibitem{debenedetti2024agentdojo}
E. Debenedetti et al.
\newblock AgentDojo: A Dynamic Environment to Evaluate Prompt Injection Attacks and Defenses.
\newblock arXiv:2406.13352, 2024.

\bibitem{zhang2024asb}
Y. Zhang et al.
\newblock Agent Security Bench: Formalizing and Benchmarking Attacks and Defenses in LLM Agents.
\newblock arXiv:2410.02644, 2024.

\bibitem{zhan2025adaptive}
Q. Zhan et al.
\newblock Adaptive Attacks Against Indirect Prompt Injection Defenses.
\newblock arXiv:2503.00061, 2025.

\bibitem{zhang2025mcpsecurity}
Y. Zhang et al.
\newblock MCP Security Bench.
\newblock arXiv:2510.15994, 2025.

\bibitem{jiang2026agentlab}
H. Jiang et al.
\newblock AgentLAB: Long-Horizon Agent Security Benchmark.
\newblock arXiv:2602.16901, 2026.

\bibitem{shi2025toolhijacker}
W. Shi et al.
\newblock ToolHijacker: Manipulating Tool Selection in LLM Agents.
\newblock arXiv:2504.19793, 2025.

\bibitem{beurerkellner2025patterns}
L. Beurer-Kellner et al.
\newblock Design Patterns for Prompt Injection Defenses.
\newblock arXiv:2506.08837, 2025.

\bibitem{yu2026toolparsing}
Z. Yu et al.
\newblock Tool Result Parsing for Robust LLM Agents.
\newblock arXiv:2601.04795, 2026.

\bibitem{wang2025agentspec}
X. Wang et al.
\newblock AgentSpec: Customizable Runtime Enforcement for LLM Agents.
\newblock arXiv:2503.18666, 2025.

\bibitem{shi2025progent}
W. Shi et al.
\newblock Progent: Privilege Control for LLM Agents.
\newblock arXiv:2504.11703, 2025.

\bibitem{kamath2025agentc}
A. Kamath et al.
\newblock Agent-C: Enforcing Temporal Constraints for LLM Agents.
\newblock arXiv:2512.23738, 2025.

\bibitem{winston2026solver}
D. Winston et al.
\newblock Solver-Aided Policy Compliance for Tool-Using Agents.
\newblock arXiv, 2026.

\bibitem{sah2026verifiertax}
A. Sah et al.
\newblock The Verifier Tax in Long-Horizon Agent Execution.
\newblock arXiv:2603.19328, 2026.

\bibitem{yao2024taubench}
S. Yao et al.
\newblock tau-bench: A Benchmark for Tool-Agent-User Interaction.
\newblock arXiv:2406.12045, 2024.

\bibitem{nandi2025sopbench}
S. Nandi et al.
\newblock SOP-Bench: Benchmarking Agents on Standard Operating Procedures.
\newblock arXiv:2506.08119, 2025.

\end{thebibliography}

\end{document}
